\documentclass[a4paper,11pt]{article}
\usepackage[scale={0.8,0.9},centering,includeheadfoot]{geometry}
\usepackage{amstext}
\usepackage{amsmath}
\usepackage{verbatim}

\begin{document}
\section*{The Gamma-distribution with constant rate}

\subsection*{Parametrisation}

The Gamma-distribution has the following density
\begin{displaymath}
    \pi(y) = \frac{b^{a}}{\Gamma(a)} y^{a-1} \exp(-by), \qquad a>0,
    \quad b>0, \quad y >0,
\end{displaymath}
where $\text{E}(y) = \mu = a/b$ and $\text{Var}(y) = 1/\tau = a/b^{2}$,
where $\tau$ is the precision and $\mu$ is the mean.

This version of the Gamma-distribution is to make the rate $b$
constant, so we use
\begin{displaymath}
    \mu = \exp(\eta)
\end{displaymath}
where $\eta$ is the linear predictor, and 
\begin{displaymath}
    \tau = (s \phi) / \mu
\end{displaymath}
where $\phi$ is the precision parameter (or $1/\phi$ is the dispersion
parameter with this definition) and $s>0$ is a fixed scaling (for the
regression model), which gives Gamma-density with
\begin{displaymath}
    a = s \phi \mu \qquad\mbox{and}\qquad b= s\phi
\end{displaymath}

\subsection*{Link-function}

The linear predictor $\eta$ is linked to the mean $\mu$ using a
default log-link
\begin{displaymath}
    \mu = \exp(\eta)
\end{displaymath}

\subsection*{Hyperparameter}

The hyperparameter is the precision parameter $\phi$, which is
represented as
\begin{displaymath}
    \phi = \exp(\theta)
\end{displaymath}
and the prior is defined on $\theta$.

\subsection*{Specification}

\begin{itemize}
\item \texttt{family="gammasv"}.
\item Required arguments: $y$ and $s$ (default value 1).
\end{itemize}

\subsubsection*{Hyperparameter spesification and default values}
%% DO NOT EDIT!
%% This file is generated automatically from models.R
\begin{description}
	\item[doc] \verb!The Gamma likelihood with constant rate!
	\item[hyper]\ 
	 \begin{description}
	 	\item[theta]\ 
	 	 \begin{description}
	 	 	\item[hyperid] \verb!58003!
	 	 	\item[name] \verb!precision parameter!
	 	 	\item[short.name] \verb!prec!
	 	 	\item[output.name] \verb!Precision-parameter for the Gammasv observations!
	 	 	\item[output.name.intern] \verb!Intern precision-parameter for the Gammasv observations!
	 	 	\item[initial] \verb!4.60517018598809!
	 	 	\item[fixed] \verb!FALSE!
	 	 	\item[prior] \verb!loggamma!
	 	 	\item[param] \verb!1 0.01!
	 	 	\item[to.theta] \verb!function(x) log(x)!
	 	 	\item[from.theta] \verb!function(x) exp(x)!
	 	 \end{description}
	 \end{description}
	\item[status] \verb!experimental!
	\item[survival] \verb!FALSE!
	\item[discrete] \verb!FALSE!
	\item[link] \verb!default log!
	\item[pdf] \verb!gammasv!
\end{description}


\subsection*{Example}

In the following example we estimate the parameters in a simulated
example.
\verbatiminput{example-gammasv.R}

\subsection*{Notes}

None.

\end{document}

% LocalWords:  hyperparameter overdispersion Hyperparameters nbinomial

%%% Local Variables: 
%%% TeX-master: t
%%% End: 
